\documentclass{situationist-preprint}

\title{Internalizing the Room}
\msid{MND-07-c}
\msfamily{MND-07}
\msdiscipline{Philosophy of Mind}
\mstype{Consequence}
\msstatus{Revised 7 October 2026}
\mskeywords{Chinese Room; internalization; systems reply; virtual mind; implementation}

\begin{document}
\maketitle

\begin{abstract}
Searle replies to the systems reply by having the man memorize the rule book and do all the work in his head. We show that, after internalization, the man executes the program and a system implemented in his brain runs it, and that these are distinct subjects. The man's failure to understand Chinese therefore says nothing about whether that system understands.
\end{abstract}

\paragraph{Note on the revision.} The version of 6 October 2026 argued that after internalization the system is identical to a part of the man's brain, and inferred that the system is not the man. The argument conflated the system with its implementation: the system is realized in the man's brain, but it is not identical to any part of it. The argument below does not rely on that identification.

\section{The argument}

Following MND-07-a, a program is \emph{executed} by an agent who carries out its operations one by one, and \emph{run} by a system whose states realize the functional organization the program specifies. After internalization, the man executes the program from memory. The system that runs it is the pattern of states and transitions realized by his carrying out the rules: a pattern whose states are, functionally, the states of a competent Chinese speaker.

\begin{proposition}\label{p:distinct}
After internalization, the system that runs the program is not the man.
\end{proposition}

\begin{proof}
The man and the system differ in their functional organization. The man's beliefs concern which characters to write in response to which; the system's states play the roles of beliefs about whatever the Chinese sentences are about. The man can be asked in English what he believes and will deny believing anything about the topics of the Chinese conversation; the system, asked in Chinese, will answer about those topics. A single subject cannot both have and lack a belief at the same time, so the two are different subjects. That both are realized in the same brain no more makes them identical than two programs running on the same computer are identical \citep{Cole1991}.
\end{proof}

By Proposition~\ref{p:distinct}, the man's lack of understanding after internalization is, as before, a fact about the executor, not about the system, and the systems reply applies unchanged \citep{Searle1980}.

\begin{thebibliography}{9}\small
\bibitem[Cole(1991)]{Cole1991} Cole, D. (1991). Artificial intelligence and personal identity. \emph{Synthese}, 88(3), 399--417.
\bibitem[Searle(1980)]{Searle1980} Searle, J.~R. (1980). Minds, brains, and programs. \emph{Behavioral and Brain Sciences}, 3(3), 417--457.
\end{thebibliography}

\end{document}
